% 本文件由 LaTeX 排版 Agent 根据有效 translation.json 自动生成。
% author=Councils; work_id=3802; title=安西拉公会议（公元314年）

\chapter{安西拉公会议（公元314年）}

% structure: p0001 source=h1 target=omitted reason=page-title-already-rendered-by-parent

马克西米努斯皇帝驾崩后不久，在加拉提亚的首府安西拉召开了一次公会议。仅有约十二位主教出席，而见于诸条规之后的签署名单并不可靠，因为其现存形式明显是后人所作，巴莱里尼弟兄已证明这一点。若我们多少可以信赖这些名单，那么似乎叙利亚和小亚细亚的几乎每个地区都有代表，因此这次会议尽管人数不多，却具有相当的分量。尚不能确定是维塔利斯（安提约基雅主教）主持，还是当时担任安西拉主教的玛尔凯鲁斯主持。\WorkTitle{会议录}将此荣誉归于后者。\par

此次公会议的纪律法令具有独特的意趣，因为它们是基督徒迫害终止后最早颁布的，并为如何妥善处理背教者作出了规定。最近发现了两份纸莎草纸文献，上面载有罗马政府发给那些背教并献祭者的官方证明。这些背教者被迫公开承认他们归属于帝国的国教，然后凭一份证明此事的文书以免去进一步的麻烦。哈纳克博士（\WorkTitle{普鲁士年鉴}）在论及背教者的屈服时写道：\par

\Quote{教会谴责此为谎言和背弃信仰，迫害结束后，这些不幸的人一部分被绝罚，一部分被迫接受严厉的纪律。谁会想到，他们耻辱的记录竟会流传到我们的时代？——然而这确实发生了。有两份这样的文书被保存下来，全然出乎意料，是埃及的沙土如此精心地保管了托付给它们的东西。第一份由克雷布斯在运往柏林的一堆纸莎草纸中发现；另一份由韦塞利在赖纳大公的纸莎草纸收藏中找到。\InnerQuote{我，狄奥革涅斯，经常献祭和奉献祭品，并在你们面前吃了祭肉，我请求你们给我一份证明。}今天，谁能不深深感动地阅读这份文书，并估量当时基督徒崩溃时所经历的苦恼和心灵恐惧呢？}\par

% structure: p0005 source=h2 target=section reason=relative-heading-rank; base=h2

\section{条规}

% structure: p0006 source=h3 target=subsection reason=relative-heading-rank; base=h2

\subsection{第一条规}

关于那些曾献祭、后来又重返斗争——并非出于伪善，而是真诚的——的司铎，已决定：他们可以保有其席位之荣誉，条件是他们未曾使用手段、安排或劝说，使自己在酷刑（仅流于表面和假装）之下显得像是受了折磨。然而，他们不得举行祭献，不得讲道，简言之，不得执行任何司铎职分的行为。\par

\EditorialNote{若望·佐纳拉斯写道：在迫害中屈服于暴君并献祭的人中，有些人在受刑后，因不能坚持到底、承受其力量和强度而失败，背弃了信仰；有些人则因懦弱，在遭受任何苦难之前就退让了，并且为了不显得是自愿献祭，他们用贿赂或恳求说服行刑者，也许对他们表现出更严厉的态度，并似乎在表面上对他们施以酷刑，以便在这种情况下献祭，从而显得是因武力而非懦弱而背弃基督。}\par

% structure: p0009 source=h3 target=subsection reason=relative-heading-rank; base=h2

\subsection{第二条规}

同样决定：那些曾献祭、后来又重返斗争的执事，可享有其他荣誉，但应远离一切神圣职务，既不可取出饼和杯，也不可作宣告。然而，若任何主教在他们身上看出心灵痛苦和谦卑的降服，主教可允予更多的宽大，或收回（已允予的）。\par

% structure: p0011 source=h3 target=subsection reason=relative-heading-rank; base=h2

\subsection{第三条规}

那些逃跑并被捉住，或被仆从出卖的人；或那些被剥夺财产、忍受酷刑、被监禁和虐待——不断宣称自己是基督徒——的人；或那些被迫接受迫害者强行塞入他们手中的东西，或接受（祭过偶像的）肉食，却持续声明自己是基督徒的人；以及那些以其全部衣着、举止和卑微的生活，始终对所发生之事表示悲伤的人：这些人，既然他们无罪，不应被拒绝于共融之外；若因某些过分的严格或无知而已经被拒绝，应立刻重新接纳。这同样适用于神职人员和教友。也已考虑：凡同样被迫而失足的教友可否被准许领受圣职；我们决定，既然他们绝无过错，可以受祝圣，条件是他们过去的生活经查是正直的。\par

\EditorialNote{在译文中，\Quote{abused}一词被用作 \GreekText{περισχισθέντας} 的对应词。佐纳拉斯将其译为 \Quote{若他们的衣服被从身上撕去}，若读法正确，这相当准确，但鲁思在博德利图书馆发现数份抄本作 \OriginalTerm{被征服}{\GreekText{περισχεθέντας}}。黑费勒采纳此读法，译为 \Quote{宣称自己是基督徒，但随后被征服，无论压迫者是否将香强塞入他们手中，或强迫他们等等}。哈蒙德译为 \Quote{并受到迫害者的骚扰，他们强行将东西塞入他们手中，或被迫等等}。无论哪种读法，措辞都是晦涩的。}\par

% structure: p0014 source=h3 target=subsection reason=relative-heading-rank; base=h2

\subsection{第四条规}

关于那些被迫献祭，并且此外还参加了偶像节庆筵席的人：凡是被拖去，但随后兴高采烈地上前，穿上最华丽的衣裳，并冷漠地享用所备筵席的，决定：这样的人应做听道者一年，跪拜者三年，仅参与祈祷两年，然后返回圆满共融。\par

\EditorialNote{希腊文中\Quote{圆满共融}一词是 \GreekText{τέλειον}（\Quote{圆满}），这是早期作者常用以指称圣体圣事的表达。参见苏伊策，\WorkTitle{词库}，该条。}\par

% structure: p0017 source=h3 target=subsection reason=relative-heading-rank; base=h2

\subsection{第五条规}

然而，凡是穿着丧服上前、坐下进食，并在整个宴席中哭泣的人，若他们已作跪拜者满三年，则应被接纳而不献祭；若他们没有进食，则应作跪拜者两年，第三年可不献祭而领受共融，以便在第四年被接纳进入圆满共融。但主教有权在考虑他们悔改的性质后，或更宽大地对待他们，或延长时间。然而，首先应彻底查考他们之前和之后的生活，据此决定宽大的程度。\par

% structure: p0019 source=h3 target=subsection reason=relative-heading-rank; base=h2

\subsection{第六条规}

关于那些仅仅是因刑罚威胁、财产充公或流放的胁迫而屈服献祭，且至今尚未悔改归正，但如今在此届会议期间带着归正意向前来的人，会议规定：他们在大日之前可作为听道者被接纳；大日之后，他们应作屈膝者三年，另两年无奉献参与共融，然后进入圆满共融，如此满六年整。若在此届会议之前已被接纳补赎，则六年应从那时起算。然而，若因病或其他原因有死亡危险或预期，则他们可被接纳，但须受限制。\par

\EditorialNote{若望·佐纳拉斯：但若那些被禁止共融的补赎者患病或因其他方式濒临死亡，他们可以被接纳共融；但须遵守这一规定或区别：若他们幸免于死并恢复健康，则应再次完全被剥夺共融，直至完成六年的补赎。}\par

% structure: p0022 source=h3 target=subsection reason=relative-heading-rank; base=h2

\subsection{第七条法令}

关于那些在外教人指定的场所参与外教宴席，但自带并食用自己食物的人，会议规定：他们作屈膝者两年后可被接纳；但是否准予奉献，须由每位主教在审查其余生行为后决定。\par

\EditorialNote{赫费勒：若干基督徒试图以世俗的精明走中间道路。一方面，希望逃避迫害，他们参与外教祭献的宴席（这些宴席在庙宇附属建筑中举行）；另一方面，为安抚良心，他们自带食物，不碰任何祭过神的东西。这些基督徒忘记了圣保禄曾命令应避免祭过神的食物，并非因为它们本身不洁（因为偶像本属虚无），而是出于另两个原因：第一，因有些人在吃时心中仍有偶像，即仍依附于偶像崇拜，从而犯罪；第二，因其他人使弟兄跌倒，从而也犯罪。此外还有第三个原因，即这些基督徒的虚伪与狡诈，他们既想显得像外教人，又想仍然作基督徒。主教会议处罚他们，使其处于第三级补赎两年，并给予每位主教权利，在此期满后，或接纳他们共融，或令其再留在第四级一段时间。}\par

% structure: p0025 source=h3 target=subsection reason=relative-heading-rank; base=h2

\subsection{第八条法令}

凡在胁迫下两次或三次献祭者，应作屈膝者四年，无奉献参与共融两年，第七年可获圆满共融。\par

% structure: p0027 source=h3 target=subsection reason=relative-heading-rank; base=h2

\subsection{第九条法令}

凡不仅背教，而且起来反对弟兄，强迫他们背教，并对此强迫行为负有责任者，应作听道者三年，再作屈膝者六年，另一年无奉献参与共融，如此满十年后，方可领受共融；在此期间，也须查考其余生行为。\par

% structure: p0029 source=h3 target=subsection reason=relative-heading-rank; base=h2

\subsection{第十条法令}

凡被立为执事者，在受祝圣时声明必须结婚，因不能持守独身，后来结了婚的，可继续担任职务，因为主教已对此予以许可。但若有对此事默不作声，在受祝圣时承诺保持现状，后来却结婚的，则应停止执事职务。\par

\Person{范·埃斯彭}评论说：\Quote{提交给主教会议并由本条例决的案件如下：当主教愿意祝圣两位执事时，其中一人声明他不打算约束自己守永久的独身，而打算结婚，因他无力保持独身；另一人则未作声。主教为每人覆手，授予执事职。}\par

祝圣后，两人都结了婚，所提出的问题是：每种情况应如何处理？主教会议裁定，那位在祝圣时作过声明的人应继续其职务，\Quote{因为主教的许可}，即他在领受执事职后可以结婚。至于保持沉默者，主教会议宣布他应停止职务。\par

主教会议对第一个问题的决议表明，存在一项普遍法律约束执事守独身；但本次主教会议认为，主教为了正当理由可以豁免此项法律，而且若主教在祝圣时（因他声明打算结婚、不能保持现状）仍祝圣了他，便视为已通过祝圣行为默示给予此许可或豁免。此外，从这一决定也可明显看出，受祝圣的执事不仅被允许结婚，而且可以在祝圣后行使婚姻。而且，那位在声明后结婚并行使婚姻的执事，不仅保留执事职，而且继续执行其职务。\par

% structure: p0034 source=h3 target=subsection reason=relative-heading-rank; base=h2

\subsection{第十一条法令}

会议规定：已订婚的贞女，若后来被他人抢走，应归还给原先与其订婚者，即使她们曾遭受强暴者的暴力。\par

\EditorialNote{参见圣巴西略致安菲洛基书第二十二条法令，其中如此规定。}\par

% structure: p0037 source=h3 target=subsection reason=relative-heading-rank; base=h2

\subsection{第十二条法令}

会议规定：凡在受洗前曾献祭，而后受了洗的，可被晋升圣职，因他们已得到洁净。\par

% structure: p0039 source=h3 target=subsection reason=relative-heading-rank; base=h2

\subsection{第十三条法令}

\OriginalTerm{乡牧主教}{Chorepiscopi}不得祝圣司铎或执事，更不得在没有主教书面委托的情况下，在另一堂区祝圣城市司铎。\par

若第十三条法令的第一部分易于理解，则第二部分却呈现出很大困难；因为城市司铎在任何情况下都无权祝圣司铎和执事，更不用说在陌生的教区。许多最有学问的人因此认为，我们所读的该法令后半部分的希腊文是不正确或有缺陷的。他们说，缺少\OriginalTerm{行某事}{\GreekText{ποιεῖν} \GreekText{τι}}，或\OriginalTerm{行某事}{aliquid agere}，即完成一项宗教礼仪。为证实此假设，他们援引了几种古代译本，尤其是\Person{伊西多尔}的译本：\Quote{但城市司铎未经主教命令，或没有其书信授权，在任何堂区（某些版本读作\OriginalTerm{在每一}{\GreekText{ἐν} \GreekText{ἐκάστῃ}}而非\OriginalTerm{在另一}{\GreekText{ἐν} \GreekText{ἑτέρᾳ}}）不可命令更多事或行事。}古代罗马法令手抄本《法令集》有相同读法，只是将\OriginalTerm{堂区}{parochia}换为\OriginalTerm{省}{provincia}。\Person{迦太基的执事富尔根提乌斯·费朗杜斯}早在很久以前就编纂了法令集，在他的《法令摘要》中同样翻译为：\Quote{城市司铎未经主教命令不得命令任何事，也不得在任何堂区行事。}\Person{范·埃斯彭}也以同样方式解释了此条法令。\par

Routh 给出了另一种解释。他声称本条典中没有缺失一词，但根据几份手抄本开头应读作 \OriginalTerm{乡村主教（与格）}{\GreekText{χωρεπισκόποις}}，下文应读作 \OriginalTerm{但肯定不}{\GreekText{ἀλλὰ} \GreekText{μὴν} \GreekText{μηδὲ}} 而非 \OriginalTerm{但不}{\GreekText{ἀλλα} \GreekText{μηδὲ}}，然后读作 \OriginalTerm{司铎（宾格）}{\GreekText{πρεσβυτέρους}} \OriginalTerm{城镇}{\GreekText{πόλεως}}，最后读作 \OriginalTerm{每个}{\GreekText{ἑκάστῃ}} 而非 \OriginalTerm{另一个}{\GreekText{ἑτέρᾳ}}，因此必须翻译为：\Quote{\Emph{乡村主教}不被允许祝圣司铎和执事（为乡村），更不用说（\OriginalTerm{但肯定不}{\GreekText{ἀλλὰ} \GreekText{μὴν} \GreekText{μηδὲ}}）未经当地主教同意，他们不能为城镇祝圣司铎。}如此根据某些手抄本（尤其是博德利图书馆的）修改的希腊文确实给出了良好的含义。不过 \OriginalTerm{但肯定不}{\GreekText{ἀλλὰ} \GreekText{μὴν} \GreekText{μηδὲ}} 并非\Quote{更不用说}之意：它的意思是\Quote{但肯定不}，这造成了相当大的差异。\par

此外，乡村主教为城镇祝圣司铎或执事的情形非常罕见；果真如此，他们已至少在本典第一部中被默示禁止。\par

% structure: p0044 source=h3 target=subsection reason=relative-heading-rank; base=h2

\subsection{第十四条典}

规定：在圣职人员中，那些戒绝肉食的司铎和执事应尝用肉，之后若愿意，可再戒绝。但若他们鄙视肉食，甚至不愿吃与肉同煮的蔬菜，且不服从本条典，则应被撤职。\par

关于希腊文读法有一场严重争议。我遵循 Routh 的意见，他依据三份手抄本、《安提约基雅若望汇编》以及拉丁译本，读作 \OriginalTerm{若他们厌恶}{\GreekText{εἰ} \GreekText{δὲ} \GreekText{βδελύσσοιντο}} 而非普通文本中的 \OriginalTerm{若他们愿意}{\GreekText{εἰ} \GreekText{δὲ} \GreekText{βούλοιντο}}，正如 Beveridge 主教此前已指出的，除非引入一个 \OriginalTerm{不}{\GreekText{μὴ}}，否则 \OriginalTerm{若他们愿意}{\GreekText{εἰ} \GreekText{δὲ} \GreekText{βούλοιντο}} 毫无意义。\par

Zonaras 指出本条典主要涉及爱筵。\par

我不同意 Hefele 对最后一句的翻译。他认为指的是\Quote{本条典}，我认为它明显指的是所谓的《宗徒典》中的第 53（52）条，\OriginalTerm{该典}{\GreekText{τῷ} \GreekText{κανόνι}} 即\Quote{那条著名的典}。\par

% structure: p0049 source=h3 target=subsection reason=relative-heading-rank; base=h2

\subsection{第十五条典}

关于在没有主教的情况下司铎可能已出售的属于教会的事物，规定：应追回教产；是否接受售价则由主教酌情决定；因为出售之物所得的收益常常可能产生更大的价值。\par

\EditorialNote{Hefele：\Quote{如果教会财产的买主通过这些财产的临时收益所获超过购买价格，主教会议认为没有必要将价格归还给他，因为他已从收益中获得充分补偿，而且根据当时有效的规则，不允许从购买款中抽取利息。此外，买主在教座空位期间购买教产是错误的。}Beveridge 和 Routh 已表明在文本中 \OriginalTerm{收回}{\GreekText{ἀνακαλεῖσθαι}}} 和 \OriginalTerm{收益}{\GreekText{πρόσοδον}} 必须读作。\par

% structure: p0052 source=h3 target=subsection reason=relative-heading-rank; base=h2

\subsection{第十六条典}

犯兽欲之罪者，若犯罪时不满二十岁，应作俯伏者十五年，然后分享祈祷；在这分享祈祷中度过五年后，即可参与献礼。但应审查他们作俯伏者的生活，然后给予宽免；若有人对其罪行贪得无厌，则应延长其俯伏期。若超过此年龄且有妻室者犯此罪，应作俯伏者二十五年，然后分享祈祷；在分享祈祷五年后，参与献礼。若超过五十岁的已婚者如此犯罪，则仅允许他们在临终时领受共融。\par

% structure: p0054 source=h3 target=subsection reason=relative-heading-rank; base=h2

\subsection{第十七条典}

与兽自污者，且患有麻风病，并传染他人者，圣主教会议命令他们应在\OriginalTerm{越冬者}{hiemantes}中祈祷。\par

% structure: p0056 source=h3 target=subsection reason=relative-heading-rank; base=h2

\subsection{第十八条典}

若有人已被祝圣为主教，但未被派往他们的教区接纳，反而侵入其他教区，迫害那里的在职主教，煽动骚乱反对他们，则应将此等之人停职并停止共融。但若他们愿意接受在司铎团中的席位，即他们先前担任司铎的地方，则不应剥夺此荣誉。但若他们对那地设立的主教继续煽动作乱，则司铎的荣誉也应被剥夺，并将他们驱逐。\par

% structure: p0058 source=h3 target=subsection reason=relative-heading-rank; base=h2

\subsection{第十九条典}

若有宣发誓愿守贞者不守其誓愿，则应使他们履行\OriginalTerm{再婚者}{digamists}的期限。此外，我们禁止守贞女子作为姐妹与男子同居。\par

% structure: p0060 source=h3 target=subsection reason=relative-heading-rank; base=h2

\subsection{第二十条典}

若任何人的妻子犯奸淫，或任何男子犯奸淫，似乎宜规定：在经过规定的补赎等级七年后，恢复其圆满共融。\par

% structure: p0062 source=h3 target=subsection reason=relative-heading-rank; base=h2

\subsection{第二十一条典}

关于犯姘居并摧毁所怀胎儿的妇女，或使用药物堕胎者，先前的规定是将她们逐出教会直到临终，有些人赞同此规定。然而，我们意欲施以稍大之宽仁，规定她们应按照规定的等级履行十年补赎。\par

% structure: p0064 source=h3 target=subsection reason=relative-heading-rank; base=h2

\subsection{第二十二条典}

关于故意杀人者，应让他们作俯伏者；但在生命终结时，可宽免他们领受圆满共融。\par

% structure: p0066 source=h3 target=subsection reason=relative-heading-rank; base=h2

\subsection{第二十三条典}

关于非故意杀人者，先前的规定指示他们按照规定的等级经七年补赎后，可被接纳领受圆满共融；但这第二条规定他们应履行五年期限。\par

% structure: p0068 source=h3 target=subsection reason=relative-heading-rank; base=h2

\subsection{第二十四条典}

从事占卜、追随异教习俗者，或带人至家中行法术或行净礼者，应受五年补赎的规定，按规定的等级：即三年作俯伏者，两年作无献礼的祈祷。\par

\EditorialNote{我读作 \OriginalTerm{外邦人的}{\GreekText{ἐθνῶν}} 代替 \OriginalTerm{时间的}{\GreekText{χρόνων}}，因而翻译为\Quote{外邦人的}。}\par

% structure: p0071 source=h3 target=subsection reason=relative-heading-rank; base=h2

\subsection{第二十五条典}

有一人已与一少女订婚，却奸污了她的妹妹，使其怀孕。之后他娶了已订婚的少女，但被奸污的少女上吊自尽。与此事相关的人被命令按照规定的等级，经十年补赎后，被接纳为\OriginalTerm{共立者}{co-standers}。\par

\SourceRecord{Translated by Henry Percival.；Nicene and Post-Nicene Fathers, Second Series；Vol. 14.；Edited by Philip Schaff and Henry Wace.；Buffalo, NY: Christian Literature Publishing Co.,；1900.；<http://www.newadvent.org/fathers/3802.htm>.}
